% ═══════════════════════════════════════════════════════════════
%  PART XI — APRIL 8, 2026: THE EIGHTH LOCK — GOLAY CODES
% ═══════════════════════════════════════════════════════════════
\clearpage
\part{Lock 8: The Universe as an Error-Correcting Code}

\section{The Extended Ternary Golay Code IS W(3,3)}
\label{sec:ternary-golay}

The extended ternary Golay code is a $[12,\,6,\,6]_3$ linear code:
length $n=k=12$, dimension $\dim = q! = 6$, minimum distance $d = q! = 6$,
over the alphabet $\GF(q) = \GF(3)$.

\begin{theorem}[Lock 8: The Ternary Golay Code]
\label{thm:lock8}
Every parameter of the extended ternary Golay code is a
W(3,3) parameter:
\begin{center}
\begin{tabular}{@{}lcl@{}}
\toprule
Code parameter & Value & W(3,3) \\
\midrule
Alphabet & $\GF(3)$ & $\GF(q)$ \\
Length & 12 & $k$ (valence) \\
Dimension & 6 & $q!$ (KO-dim) \\
Distance & 6 & $q!$ \\
Codewords & 729 & $q^{q!} = 3^6$ \\
Rate & 1/2 & self-dual \\
$A_k = A_{12}$ & 24 & $f$ (lines) \\
\bottomrule
\end{tabular}
\end{center}
The weight-$q!$ codewords form the Steiner system $S(5,6,12)$
with $132$ hexads, and the automorphism group is $2\cdot M_{12}$.
\end{theorem}

The number of maximum-weight codewords, $A_{12} = 24 = f$,
equals the number of lines in the W(3,3) design.
The Mathieu group $M_{12}$ has order
$|M_{12}| = k!/(k{-}5)! = 12\cdot 11\cdot 10\cdot 9\cdot 8 = 95040$
and is \textbf{sharply 5-transitive} on $k=12$ points.


\section{The Binary Golay Code: The Bosonic Shadow}
\label{sec:binary-golay}

The extended binary Golay code is a $[24,\,12,\,8]_2$ code:
\begin{equation}
[f,\; k,\; 2^q]_2 = [24,\; 12,\; 8]_2.
\end{equation}
Its automorphism group is $M_{24}$, which leads to the
Conway group, the Leech lattice $\Lambda_{24}$, and ultimately
the Monster group.

\begin{theorem}[Two Golay Codes $=$ Two Faces of W(3,3)]
\begin{center}
\begin{tabular}{@{}lcccc@{}}
\toprule
& Alphabet & Length & Dim & Distance \\
\midrule
Ternary Golay & $\GF(q)$ & $k=12$ & $q!=6$ & $q!=6$ \\
Binary Golay & $\GF(\lambda)$ & $f=24$ & $k=12$ & $2^q=8$ \\
\bottomrule
\end{tabular}
\end{center}
The ternary code encodes the fermionic sector
(alphabet $=$ field order, length $=$ valence).
The binary code encodes the bosonic sector
(length $=$ lines, dimension $=$ valence).
Together they generate the complete moonshine tower:
$M_{12}\subset M_{24}\to\mathrm{Co}_1\to\Lambda_{24}\to\mathbb{M}$.
\end{theorem}


\section{The Division Algebra Connection}
\label{sec:furey}

Furey~\cite{Furey2022} demonstrates that one Standard Model generation
is encoded in $\mathbb{R}\otimes\mathbb{C}\otimes\mathbb{H}\otimes\mathbb{O}$,
whose complex dimension is
\begin{equation}
1\times 2\times 4\times 8 = 64 = 2^{q!} = \mu^q.
\end{equation}
This is simultaneously the dimension of $\Cl(q!)$, the number of
irreducible representations of $\Cl_q(q,\lambda)$, and the total
number of sub-faces of the $(q!{-}1)$-simplex (Pascal row~$q!$).
One chiral generation has $32 = 2^{q+\lambda}$ Weyl fermions.


\section{The Information-Theoretic Interpretation}
\label{sec:info-theory}

The complete dictionary now reads:
\begin{quote}
\emph{The universe is an error-correcting code over $\GF(q)=\GF(3)$.
The code length is $k=12$ (the W(3,3) valence),
the dimension is $q!=6$ (the KO-dimension of the SM),
and the minimum distance is $q!=6$ (optimal error protection).
The code is self-dual ($C=C^\perp$, mirroring matter--antimatter
symmetry), unique (Golay), and perfect (optimal sphere-packing
in Hamming space).}

\emph{The symmetry group $M_{12}$ is a sporadic simple group,
sharply 5-transitive on $k$ points, embedding into $M_{24}$
and thence into the Monster via the Leech lattice.
This is the eighth lock: the ternary Golay code,
living over $\GF(q)$ with parameters $(k, q!, q!)$,
is one more independent mathematical structure that
uniquely requires $q=3$.}
\end{quote}

\begin{thebibliography}{9}
\bibitem{Furey2022} N. Furey and M.\,J. Hughes,
  ``One generation of Standard Model Weyl representations as a
  single copy of $\mathbb{R}\otimes\mathbb{C}\otimes\mathbb{H}
  \otimes\mathbb{O}$,'' \emph{Phys.\ Lett.\ B} \textbf{831} (2022) 136959.
\end{thebibliography}
